
\documentclass[11pt]{article}
\usepackage[margin=1in]{geometry}
\usepackage{amsmath,amssymb,booktabs,hyperref,graphicx,parskip,enumitem}
\usepackage{xcolor}
\hypersetup{colorlinks=true,linkcolor=blue!50!black,urlcolor=blue!50!black}
\setlist{nosep}

\title{Replication Report: (title not recorded)\\[2pt]\large OSTI 1427646 \textbar{} Coverage 8/10, Agreement 8/10}
\author{Rick Stevens \and Ollie (AI Assistant)\\\small Argonne National Laboratory}
\date{2026-04-27}

\begin{document}\maketitle

\section{Source paper}
\begin{itemize}
\item \textbf{Title:} (title not recorded)
\item \textbf{Authors:} Not recorded
\item \textbf{Year:} n/a
\item \textbf{Domain:} Not recorded
\item \textbf{Identifier:} OSTI 1427646
\end{itemize}


\section{Replication objective}
The central claim of the paper concerns not recorded. Our objective was to reproduce the headline numerical/qualitative results within the constraints of the AI-driven Replication Project (24--72\,h, commodity GPU/CPU compute, no author contact). A replication plan is archived at \texttt{1427646-Deep-Learning-of-Atomically-Resolved-Scanning/replication\_plan.pdf}.

\section{Methodology}
We followed the standard Replication Project workflow: ingest the source PDF, extract method and data pointers, draft a replication plan, attempt the smallest end-to-end pipeline that produces a comparable number, then scale up where compute and data permit. Tools used in this replication are catalogued in the meta-paper's computational tools inventory (see \texttt{COMPUTATIONAL\_TOOLS\_INVENTORY.md}).

This paper went through the Tier-Lift pipeline; supplementary notes at \texttt{.openclaw/workspace/24h-progress/batch3-materials/report\_1427646\_STEM.md}, \texttt{.openclaw/workspace/24h-progress/tier-lift-v2.5/1427646/REPORT.md}, \texttt{.openclaw/workspace/24h-progress/tier-lift-v2/1427646.md}.

\subsection*{Paper-directory README excerpt}
\small
\par\medskip\textbf{Deep Learning of Atomically Resolved Scanning Transmission Electron Microscopy Images: Chemical Identification and Tracking Local Transformations}\par
\par$\bullet$\ **OSTI ID:** 1427646
\par$\bullet$\ **Rank:** \#10
\par$\bullet$\ **Replication Score:** 9/10 (plan) — achieved **7/10** with synthetic physics surrogate
\par$\bullet$\ **Replication date:** 2026-04-24
\par$\bullet$\ **Compute:** uicgpu, NVIDIA A100 80GB (GPU 4)
\par$\bullet$\ **Wall time:** \textasciitilde{}2 min end-to-end (data gen + 30 epochs + eval + figures)

\par\medskip\textbf{\# Why This Paper}\par
Structured STEM image datasets + deep learning for quantitative atomic classification — a clear computational pipeline suitable for AI-driven replication.

\par\medskip\textbf{\# What the paper does}\par
Trains a fully convolutional network on multislice-simulated HAADF-STEM images of SrTiO3 / La0.7Sr0.3MnO3 [001] heterostructures; applies it to experimental STEM frames to produce per-pixel chemical identity maps (Sr, Ti, La/Sr, Mn, vacuum) and track antisite / vacancy transformations.

\par\medskip\textbf{\# What we replicated}\par
\par$\bullet$\ **Model:** 7.76M-param U-Net (5 encoder stages, BN+ReLU, skip connections, softmax head) — PyTorch 2.5.1 / CUDA 12.1.
\par$\bullet$\ **Training protocol:** AdamW, cosine LR, cross-entropy with class weighting, on-the-fly flip/rotate/brightness aug — 30 epochs, batch 16.
\par$\bullet$\ **Data:** Physics-motivated synthetic surrogate (not full multislice) — perovskite [001] square lattice with A/B sublattice offsets, Z\textasciicircum{}1.7 HAADF-like intensity, random interface diffuseness, Poisson noise, scan-line jitter, A-site/B-site vacancies + antisite defects. 512 train / 64 val / 128 test frames (256×256).
\par$\bullet$\ **Evaluation:** Pixel-level accuracy + per-class precision/recall/F1; peak detection vs. ground-truth column centres (RMSE + det. precision/recall).

\par\medskip\textbf{\# Results (held-out test, 128 frames, 8.4M pixels)}\par
\par$\bullet$\ Pixel accuracy: **0.988**
\par$\bullet$\ Per-class F1: Sr 0.994, Ti 0.949, La/Sr 0.998, Mn 0.958
\par$\bullet$\ Column-centre RMSE (at \textasciitilde{}0.2 Å/px): Sr 0.11 px, Ti 0.47 px, La/Sr 0.06 px, Mn 0.45 px — sub-pixel for all classes
\par$\bullet$\ Detection recall ≥ 0.98 for every atom class

Consistent with the paper's qualitative claim that an FCN trained purely on simulation produces quantitatively useful chemical mapping of perovskite interfaces.

\par\medskip\textbf{\# Status}\par
\par$\bullet$\ [x] Paper reviewed
\par$\bullet$\ [x] Plan identified (plan PDF in this directory)
\par$\bullet$\ [x] Code implemented (`replication/src/`)
\par$\bullet$\ [x] Trained on A100 (uicgpu)
\par$\bullet$\ [x] Results validated on synthetic test set
\par$\bullet$\ [x] Report written (`replication/report/report.pdf`)

[\textbackslash{}textit\{README truncated\}]
\normalsize

The replication was driven by an OpenClaw agent loop (Claude Opus / GPT-5 class models) issuing shell, Python, and (where applicable) GPU jobs. Compute usage and wall-time for individual papers are summarised in the meta-paper's resource table; not separately recorded for this paper unless noted in the Tier-Lift artefact. Sub-agents were spawned for replication-plan generation and (for papers in batches 1--4) for tier-lift extension runs.

\section{What was reproduced}
\begin{itemize}
\item Not recorded.
\end{itemize}

\section{What was missing or incomplete}
\begin{itemize}
\item Not recorded.
\end{itemize}
The dominant blocker categories for this paper, mapped onto the meta-paper's 9-category friction taxonomy (\S4), are listed in \S\ref{sec:friction} below.

\section{Quantitative agreement}
Per-quantity numerical comparisons are not separately tabulated in the structured evaluation record for this paper beyond the agreement rationale below; where the rationale references specific numbers, they are reproduced verbatim. A full numerical comparison table is recorded in the per-replication artefacts (see \S\ref{sec:repro}). For papers below 8/8, the agreement rationale identifies the specific quantities that diverge.

\paragraph{Agreement rationale.} Not recorded.

\section{Score and friction-point tagging}\label{sec:friction}
\begin{itemize}
\item \textbf{Coverage:} 8/10 --- Not recorded.
\item \textbf{Agreement:} 8/10 --- Not recorded.
\end{itemize}

\noindent\textbf{Friction points encountered (taxonomy tags):}
\begin{itemize}
\item None recorded.
\end{itemize}

\section{Follow-on questions}
\begin{itemize}
\item Not recorded.
\end{itemize}

\section{Reproducibility pointers}\label{sec:repro}
\begin{itemize}
\item \textbf{Paper directory:} \texttt{1427646-Deep-Learning-of-Atomically-Resolved-Scanning}
\item \textbf{Replication artefacts:}
\begin{itemize}
\item \texttt{/Users/stevens/Dropbox/REPLICATE-PROJECT/1427646-Deep-Learning-of-Atomically-Resolved-Scanning/replication}
\end{itemize}
\item \textbf{Replication-dir hint from JSONL:} \texttt{/Users/stevens/Dropbox/REPLICATE-PROJECT/1427646-Deep-Learning-of-Atomically-Resolved-Scanning}
\item \textbf{Status:} tier\_lift\_v2.5\_2026\_04\_27
\end{itemize}

To re-run, change directory into the paper directory above and consult its \texttt{README.md} for the entry-point script. Most replications follow the convention \texttt{replication/run\_*.sh} or \texttt{replication/<subdir>/run.py}.

\end{document}
